\documentclass[10pt,a4paper]{amsart}
\usepackage{amsmath,amssymb,amsthm,graphicx}
\usepackage{verbatim}
\newenvironment{reference}{\comment}{\endcomment}
%\newenvironment{reference}{}{}
\newenvironment{slogan}{\comment}{\endcomment}
\newenvironment{history}{\comment}{\endcomment}

% For commutative diagrams we use Xy-pic
\usepackage[all]{xy}

% We use 2cell for 2-commutative diagrams.
\xyoption{2cell}
% Theorem environments.
%
\theoremstyle{plain}
\newtheorem{theorem}[subsection]{Theorem}
\newtheorem{proposition}[subsection]{Proposition}
\newtheorem{lemma}[subsection]{Lemma}

\theoremstyle{definition}
\newtheorem{definition}[subsection]{Definition}
\newtheorem{example}[subsection]{Example}
\newtheorem{exercise}[subsection]{Exercise}
\newtheorem{situation}[subsection]{Situation}

\theoremstyle{remark}
\newtheorem{remark}[subsection]{Remark}
\newtheorem{remarks}[subsection]{Remarks}

\numberwithin{equation}{subsection}

% Macros
%
\def\lim{\mathop{\mathrm{lim}}\nolimits}
\def\colim{\mathop{\mathrm{colim}}\nolimits}
\def\Spec{\mathop{\mathrm{Spec}}}
\def\Hom{\mathop{\mathrm{Hom}}\nolimits}
\def\Ext{\mathop{\mathrm{Ext}}\nolimits}
\def\SheafHom{\mathop{\mathcal{H}\!\mathit{om}}\nolimits}
\def\SheafExt{\mathop{\mathcal{E}\!\mathit{xt}}\nolimits}
\def\Sch{\mathit{Sch}}
\def\Mor{\mathop{\mathrm{Mor}}\nolimits}
\def\Ob{\mathop{\mathrm{Ob}}\nolimits}
\def\Sh{\mathop{\mathit{Sh}}\nolimits}
\def\NL{\mathop{N\!L}\nolimits}
\def\CH{\mathop{\mathrm{CH}}\nolimits}
\def\proetale{{pro\text{-}\acute{e}tale}}
\def\etale{{\acute{e}tale}}
\def\QCoh{\mathit{QCoh}}
\def\Ker{\mathop{\mathrm{Ker}}}
\def\Im{\mathop{\mathrm{Im}}}
\def\Coker{\mathop{\mathrm{Coker}}}
\def\Coim{\mathop{\mathrm{Coim}}}

% Boxtimes
%
\DeclareMathSymbol{\boxtimes}{\mathbin}{AMSa}{"02}

%
% Macros for moduli stacks/spaces
%
\def\QCohstack{\mathcal{QC}\!\mathit{oh}}
\def\Cohstack{\mathcal{C}\!\mathit{oh}}
\def\Spacesstack{\mathcal{S}\!\mathit{paces}}
\def\Quotfunctor{\mathrm{Quot}}
\def\Hilbfunctor{\mathrm{Hilb}}
\def\Curvesstack{\mathcal{C}\!\mathit{urves}}
\def\Polarizedstack{\mathcal{P}\!\mathit{olarized}}
\def\Complexesstack{\mathcal{C}\!\mathit{omplexes}}
% \Pic is the operator that assigns to X its picard group, usage \Pic(X)
% \Picardstack_{X/B} denotes the Picard stack of X over B
% \Picardfunctor_{X/B} denotes the Picard functor of X over B
\def\Pic{\mathop{\mathrm{Pic}}\nolimits}
\def\Picardstack{\mathcal{P}\!\mathit{ic}}
\def\Picardfunctor{\mathrm{Pic}}
\def\Deformationcategory{\mathcal{D}\!\mathit{ef}}
\usepackage{cleveref}
\crefname{originallemma}{original lemma}{original lemmas}
\Crefname{originallemma}{Original lemma}{Original lemmas}
\makeatletter
\newlabel{lemma-stalk-push-closed}{{0DK1}{source}{}{sourcepointer-1}{}}
\newlabel{lemma-stalk-push-closed@cref}{{[originallemma][0][]0DK1}{[1][0][]source}{}{}{}}
\makeatother
\crefname{originallemma}{original lemma}{original lemmas}
\Crefname{originallemma}{Original lemma}{Original lemmas}
\makeatletter
\newlabel{spaces-properties-lemma-stalk-pullback}{{04K2}{source}{}{sourcepointer-2}{}}
\newlabel{spaces-properties-lemma-stalk-pullback@cref}{{[originallemma][0][]04K2}{[1][0][]source}{}{}{}}
\makeatother
\crefname{originaltheorem}{original theorem}{original theorems}
\Crefname{originaltheorem}{Original theorem}{Original theorems}
\makeatletter
\newlabel{spaces-properties-theorem-exactness-stalks}{{04K5}{source}{}{sourcepointer-3}{}}
\newlabel{spaces-properties-theorem-exactness-stalks@cref}{{[originaltheorem][0][]04K5}{[1][0][]source}{}{}{}}
\makeatother
\title{The projection formula for a closed immersion of algebraic spaces}
\author{The Stacks Project Authors}
\date{stacks-project:spaces-morphisms.tex:local-snapshot-20261012; GNU FDL 1.2 or later}
\begin{document}
\maketitle
\setcounter{subsection}{1}
\noindent\textit{Editorial note.} All delivered mathematical text is the author's own, copied from the pinned Stacks Project chapter file spaces-morphisms.tex; no statement or proof step was rewritten, reordered or removed.\par\medskip
\noindent\textit{Editorial note.} The statement involves a morphism of ringed topoi induced by the closed immersion; that construction is part of the statement span and the proof uses only the two adjunction maps written there.\par\medskip
\noindent\textit{Editorial note.} The three references of the counted spans are recorded as non-load-bearing navigation pointers; the argument that consumes them stays inside the counted proof span.\par\medskip
\noindent\textit{Editorial note.} Preamble: the three native excerpt spans of sources/stacks-spaces-morphisms/preamble.tex (lines 8-13, 14-18 and 154-217). The chapter file itself carries no preamble text.\par\medskip
\noindent\textit{Allowed standard prerequisites (editorial).}\begin{itemize}
\item Exactness of the stalk functor on the etale site of an algebraic space. Properties of Spaces, Theorem spaces-properties-theorem-exactness-stalks, used at proof lines 2672-2674. Reduces the isomorphism of sheaves to an isomorphism on stalks.
\item Behaviour of stalks under pullback of sheaves. Properties of Spaces, Lemma spaces-properties-lemma-stalk-pullback, used at proof lines 2688-2691. Identifies the stalks of the pullbacks appearing in the computation.
\item Behaviour of stalks under pushforward along a closed immersion. Lemma lemma-stalk-push-closed, used at proof lines 2691-2693. Identifies the stalk of f\_*G at xbar with the stalk of G at zbar.
\item The projection formula as an adjunction between pullback and pushforward. used at proof lines 2662-2671. Constructs the canonical map whose isomorphism is asserted.
\end{itemize}\medskip
\ifdefined\hypertarget\hypertarget{sourcepointer-1}{}\fi
\noindent\textit{Source reference.} Original lemma 0DK1: focused-seeds/AGP/sources/stacks-snapshot-20261011/tags.txt, source lines 16996--16996. Navigation pointer into the Stacks Project; the standard fact it supplies is applied inside the counted proof span and the argument resting on it stays there.\par\medskip
\ifdefined\hypertarget\hypertarget{sourcepointer-2}{}\fi
\noindent\textit{Source reference.} Original lemma 04K2: focused-seeds/AGP/sources/stacks-snapshot-20261011/tags.txt, source lines 5958--5958. Navigation pointer into the Stacks Project; the standard fact it supplies is applied inside the counted proof span and the argument resting on it stays there.\par\medskip
\ifdefined\hypertarget\hypertarget{sourcepointer-3}{}\fi
\noindent\textit{Source reference.} Original theorem 04K5: focused-seeds/AGP/sources/stacks-snapshot-20261011/tags.txt, source lines 5961--5961. Navigation pointer into the Stacks Project; the standard fact it supplies is applied inside the counted proof span and the argument resting on it stays there.\par\medskip
\begin{lemma}
\label{lemma-closed-immersion-rings}
Let $S$ be a scheme. Let $i : Z \to X$ be a closed immersion of algebraic
spaces over $S$. Let $\mathcal{A}$ be a sheaf of rings on $X_\etale$.
Let $\mathcal{B}$ be a sheaf of rings on $Z_\etale$.
Let $\varphi : \mathcal{A} \to i_{small, *}\mathcal{B}$
be a homomorphism of sheaves of rings so that we obtain a
morphism of ringed topoi
$$
f : (\Sh(Z_\etale), \mathcal{B}) \longrightarrow (\Sh(X_\etale), \mathcal{A}).
$$
For a sheaf of $\mathcal{A}$-modules $\mathcal{F}$ and a
sheaf of $\mathcal{B}$-modules $\mathcal{G}$ the canonical map
$$
\mathcal{F} \otimes_\mathcal{A} f_*\mathcal{G}
\longrightarrow
f_*(f^*\mathcal{F} \otimes_\mathcal{B} \mathcal{G}).
$$
is an isomorphism.
\end{lemma}

\begin{proof}
The map is the map adjoint to the map
$$
f^*\mathcal{F} \otimes_\mathcal{B}
f^* f_*\mathcal{G} =
f^*(\mathcal{F} \otimes_\mathcal{A} f_*\mathcal{G})
\longrightarrow
f^*\mathcal{F} \otimes_\mathcal{B} \mathcal{G}
$$
coming from $\text{id} : f^*\mathcal{F} \to f^*\mathcal{F}$
and the adjunction map $f^* f_*\mathcal{G} \to \mathcal{G}$.
To see this map is an isomorphism, we may check on stalks
(Properties of Spaces, Theorem
\ref{spaces-properties-theorem-exactness-stalks}).
Let $\overline{z} : \Spec(k) \to Z$ be a geometric point with
image $\overline{x} = i \circ \overline{z} : \Spec(k) \to X$.
Working out what our maps does on stalks, we see that we
have to show
$$
\mathcal{F}_{\overline{x}}
\otimes_{\mathcal{A}_{\overline{x}}}
\mathcal{G}_{\overline{z}} =
(\mathcal{F}_{\overline{x}}
\otimes_{\mathcal{A}_{\overline{x}}}
\mathcal{B}_{\overline{z}}) \otimes_{\mathcal{B}_{\overline{z}}}
\mathcal{G}_{\overline{z}}
$$
which holds true. Here we have used that
taking tensor products commutes with taking stalks, the
behaviour of stalks under pullback
Properties of Spaces, Lemma \ref{spaces-properties-lemma-stalk-pullback}, and
the behaviour of stalks under pushforward along a closed immersion
Lemma \ref{lemma-stalk-push-closed}.
\end{proof}

\end{document}
